本节第 9 节报告的全部计算是在 CERN 的一台专用 PC 机群上完成的。
作者感谢 CERN 管理层提供所需经费，
并感谢 CERN IT 部门的技术支持。



\lusapp A. 记号约定

\vskip-2.5ex

\lussubsec A.1 规范群

$\SUn$ 的李代数可以等同于全部复的、反厄米、无迹 $N\times N$ 矩阵构成的线性空间。
取这类矩阵的一组基 $T^a$，$a=1,\ldots,N^2-1$，
李代数的一般元素 $X$ 即 $X=X^aT^a$，分量 $X^a$ 为实数
（重复的群指标自动求和）。
假定生成元满足
\equation{
  \tr\{T^aT^b\}=-\frac{1}{2}\delta^{ab},
  \enum
}
除此之外不作进一步规定。

规范场的归一化使得规范变换与协变导数
不显含规范耦合。
$D$ 维洛伦兹指标 $\mu,\nu,\ldots$ 取值 $0$ 到 $D-1$，
成对出现时求和。
时空度规取欧几里得度规。特别地，
对任意动量 $p$ 有 $p^2=p_{\mu}p_{\mu}$ 且 $\delta_{\mu\mu}=D$。


\lussubsec A.2 狄拉克矩阵

狄拉克矩阵 $\dirac{\mu}$ 满足
\equation{
  \diracstar{\mu}{\dagger}=\dirac{\mu},\qquad
  \{\dirac{\mu},\dirac{\nu}\}=2\delta_{\mu\nu}.
  \enum
}
标量积 $\dirac{\mu}p_{\mu}$ 照 usual 缩写为
$\slash{p}$，且
$\sigma_{\mu\nu}={i\over2}\left[\dirac{\mu},\dirac{\nu}\right]$。
在任何接近 4 的维数 $D$ 下，
\equation{
  \dirac{5}=\dirac{0}\dirac{1}\dirac{2}\dirac{3},
  \qquad
  \tr\{1\}=4.
  \enum
}
注意：若按这些约定使用维数正规化，
则轴矢流需要重正化。


\lussubsec A.3 格点理论

本文考虑的格点理论照 usual 建立在间距 $a=1$ 的四维超立方格点上。
夸克场采用 Wilson 表述：
夸克场位于格点座上，
带有与连续理论中夸克场相同的指标。
在规范场 $U(x,\mu)$ 存在下，
作用于夸克场 $\psi(x)$ 的规范协变前向与后向差分算符定义为
\equation{
  \nab{\mu}\psi(x)=U(x,\mu)\psi(x+\hat{\mu})-\psi(x),
  \enum
  \nexteq{2.5ex}
  \nabstar{\mu}\psi(x)=\psi(x)-U(x-\hat{\mu},\mu)^{-1}\psi(x-\hat{\mu}),
  \enum
}
其中 $\hat{\mu}$ 表示 $\mu$ 方向的单位矢量。
在格点上，$\drv{\mu}$ 与 $\drvstar{\mu}$ 表示普通的前向与后向差分算符，
而
\equation{
  \ring{\partial}_{\mu}=\frac{1}{2}(\drv{\mu}+\drvstar{\mu})
  \enum
}
表示它们的对称组合。

微分算符 $\partial^a_{x,\mu}$ 按下式作用于规范场 $U$ 的可微函数 $f(U)$：
\equation{
  \partial^a_{x,\mu}f(U)=
  {\rmd\over\rmd s}f(\rme^{sX}U)\!\left.{{\vphantom{r\over s}}}\right|_{s=0},
  \hskip1.7em
  X(y,\nu)=\cases{T^a & if $(y,\nu)=(x,\mu)$,\cr
                          \noalign{\vskip0.8ex}
                        0  & otherwise.\cr}
  \enum
}
虽然这些算符依赖于生成元 $T^a$ 的选取，
可以证明梯度场
\equation{
  (\partial f)(U,x,\mu)=T^a\partial^a_{x,\mu}f(U)
  \enum
}
与基的选取无关。



\lusapp B. 映射 (5.9) 的性质

映射 (5.9) 可微且满足
\equation{
  \Psu(\rme^{sX})\mathrel{\mathop=_{s\to0}}
  sX+\rmO(s^2)
  \enum
}
对所有属于 $\SUn$ 李代数的矩阵 $X$ 成立。
这些性质意味着存在 $\SUn$ 中单位元的一个开邻域，
$\Psu$ 在其中是一一对应的，
且雅可比矩阵
\equation{
  J(U)^{ab}=-2\kern1pt\tr\kern-1pt\left\{T^a\Psu(T^bU)\right\}
  \enum
}
处处非奇异。

利用简单的范数估计可以证明：
区域
\equation{
  \Dsu=\left\{U\in\SUn\bigm|
  \Re\tr\{1-U\}\leq(16N)^{-1}\right\}
  \enum
}
的内部就是这样一个单位元的邻域。
此外，雅可比行列式 $\det J(U)$ 在 $\Dsu$ 上严格为正。
于是，对该 $\SUn$ 不变测度（在李代数上）
相应的狄拉克 $\delta$ 函数而言，
恒等式
\equation{
  \int_{\Dsu} \rmd U\,\delta\left(\Psu(U)-\Psu(V)\right)f(U)=
  k\det J(V)^{-1}f(V)
  \enum
}
对所有可积函数 $f(U)$ 及所有 $V\in\Dsu$ 成立。
此式中的比例常数 $k$ 为正且可以显式求出，
但本文不需要它的数值。



\lusapp C. Wick 收缩

O($a$) 改进格点理论中基本费米子场的 Wick 收缩，
先就时间步长为 $\eps$ 的离散化流方程情形算出。
这样可以避免时序上的歧义，
而在计算的最后容易取极限 $\eps\to0$。


\lussubsec C.1 费米子作用量与泛函积分

作为开始，回顾一下是有帮助的：
$4+1$ 维 O($a$) 改进格点理论的 total 作用量的费米子部分是
一个边界项
\equation{
  \SF=\sum_x\bigl\{\psibar(x)\Dm\psi(x)
  +\cfl\lambdabar(0,x)\lambda(0,x)\bigr\},
  \enum
}
与体作用量 (5.8) 之和。
式 (C.1) 中的 Wilson--Dirac 算符 $\Dm$
包含 Sheikholeslami--Wohlert 项 [\cite{SW},\cite{SFimp}]
以及带裸质量矩阵 $M_0$ 的质量项。

在体作用量 (5.8) 中，
随时间演化场理解为满足边界条件 (2.6)。
于是泛函积分中被积分的、不受约束的费米子场
（反费米子场类似）可取为
\equation{
  \psi(x),\quad\chi(\eps,x),\ldots,\chi(T,x),\quad
  \lambda(0,x),\ldots,\lambda(T-\eps,x).
  \enum
}
由于作用量对这些场是二次的，
费米子积分由 Wick 定理给出。


\lussubsec C.2 场变换

在下文中，满足
\equation{
  K_{\eps}(t,x;t,y)=\delta_{xy}
  \enum
}
（在范围 $0\leq t\leq T$ 内的所有 $t$ 处）
以及离散化夸克流方程
\equation{
  \left(\partial_t-\Delta\right)K_{\eps}(t,x;s,y)=0
  \enum
}
（若 $0\leq s\leq t<T$）
的核 $K_{\eps}(t,x;s,y)$ 扮演着重要角色。
显然，由于流方程 (C.4) 可以从时刻 $t=s$
开始一步步求解，
这些方程在所有 $t\geq s$ 处唯一确定该核。

通过一个场变换可以简化费米子传播子的计算：
将 $\chi(t,x)$ 通过 $\psi(x)$ 与另一个定义在
$t=\eps,\ldots,T$ 的场 $\phi(t,x)$ 表示为
\equation{
  \chi(t,x)=\sum_yK_{\eps}(t,x;0,y)\psi(y)
  +\eps\sum_{s=\eps}^t\sum_yK_{\eps}(t,x;s,y)\phi(s,y).
  \enum
}
对任意给定的场 $\psi(x)$，
由场变量 $\chi(\eps,x),\ldots,\chi(T,x)$
到新变量 $\phi(\eps,x),\ldots,\phi(T,x)$ 的映射
由于恒等式
\equation{
  \phi(t+\eps,x)=\cases{
  \left(\partial_t-\Delta\right)\chi(t,x) & if $0<t<T$,\cr
  \noalign{\vskip1.5ex}
  {1\over\eps}(\chi(\eps,x)-\psi(x))-\Delta\psi(x) &
  if $t=0$.\cr
  }
  \enum
}
而是可逆的。
而且相应的雅可比矩阵是下三角的，
其行列式不依赖规范场。

经过该场变换以及反费米子场的相应变换后，
费米子作用量取如下形式：
\equation{
  \SF+\SFfl=\sum_x
  \bigl\{\psibar(x)\Dm\psi(x)+\cfl\lambdabar(0,x)\lambda(0,x)\bigr\}
  \noenum
  \nexteq{0.0ex}
  {\phantom{\SF+\SFfl=}}
  +\eps\sum_{t=0}^{T-\eps}\sum_x\bigl\{
  \lambdabar(t,x)\phi(t+\eps,x)+\phibar(t+\eps,x)\lambda(t,x)\bigr\}.
  \enum
}
这一变换从而实现了场变量之间近乎完全的退耦。


\lussubsec C.3 收缩

由于基本夸克场 $\psi,\psibar$ 与其他场完全退耦，
这些场的收缩
\equation{
  \wick{\psi(x)\psibar(y)}=S(x,y),
  \enum
  \nexteq{2.5ex}
  \Dm S(x,y)=\delta_{xy},
  \enum
}
如预期那样与四维理论中的夸克传播子一致。
由作用量 (C.7) 导出的其余非零收缩为
\equation{
  \verylongwick{\phi(t+\eps,x)\lambdabar(t,y)}=
  \longwick{\lambda(t,x)\phibar(t+\eps,y)}={1\over\eps}\delta_{xy},
  \quad t=0,\eps,\ldots T-\eps,
  \enum
  \nexteq{2.5ex}
  \longwick{\phi(\eps,x)\phibar(\eps,y)}=-{\cfl\over\eps^2}\delta_{xy}.
  \enum
}
对于涉及场 $\chi$ 与 $\chibar$ 的收缩则得到
\equation{
  \longwick{\chi(t,x)\lambdabar(s,y)}=
  \theta(t>s)K_{\eps}(t,x;s+\eps,y),
  \enum
  \nexteq{3.0ex}
  \longwick{\lambda(t,x)\chibar(s,y)}=
  \theta(t<s)K_{\eps}(s,y;t+\eps,x)^{\dagger},
  \enum
  \nexteq{3.0ex}
  \longwick{\chi(t,x)\psibar(y)}=\sum_{v}
  K_{\eps}(t,x;0,v)S(v,y),
  \enum
  \nexteq{2.4ex}
  \wick{\psi(x)\chibar(s,y)}=\sum_{w}
  S(x,w)K_{\eps}(s,y;0,w)^{\dagger},
  \enum
  \nexteq{2.4ex}
  \longwick{\chi(t,x)\chibar(s,y)}=\sum_{v,w}
  K_{\eps}(t,x;0,v)S(v,w)
  K_{\eps}(s,y;0,w)^{\dagger}
  \noenum
  \nexteq{1.5ex}
  {\phantom{\longwick{\chi(t,x)\chibar(s,y)}=}}
  -\cfl\sum_vK_{\eps}(t,x;\eps,v)K_{\eps}(s,y;\eps,v)^{\dagger}.
  \enum  
}
在所有这些等式中，
$\chi$ 与 $\chibar$ 的流时间变量取值于 $(0,T]$，
而 $\lambda$ 与 $\lambdabar$ 的流时间变量取值于 $[0,T)$。


\lussubsec C.4 连续时间极限下的收缩

当时间步长 $\eps$ 趋于零时，
核 $K_{\eps}(t,x;s,y)$ 光滑地收敛到格点理论连续时间表述下
夸克流方程的基本解 $K(t,x;s,y)$。
因此，由式 (C.8) 与 (C.12)--(C.16) 给出的
费米子场全部非零收缩的集合具有良定义的极限。



\lusapp D. 格点流方程的积分

纯规范梯度流的数值积分已在文献~[\cite{WilsonFlow}] 的附录 C 中讨论过。
这里把该文献中描述的三阶龙格--库塔积分器推广到夸克流方程。


\lussubsec D.1 梯度流的积分

仿照文献~[\cite{WilsonFlow}]，
把流方程 (5.2) 写成抽象形式
\equation{
   \partial_tV_t=Z(V_t)V_t,
   \enum
}
其中规范场 $V_t$
被视为一个（高维）李群的元素，
而 Wilson 作用量的梯度 $Z(V_t)$ 是相应李代数的元素。
上述龙格--库塔积分器以步长 $\eps$ 推进。
设 $V_t$ 在某流时间 $t$ 处已知，
通过计算场
\equation{
  W_0=V_t,
  \noenum
  \nexteq{2.5ex}
  W_1=\exp\bigl\{\frac{1}{4}Z_0\bigr\}W_0,
  \noenum
  \nexteq{2.5ex}
  W_2=\exp\bigl\{\frac{8}{9}Z_1-\frac{17}{36}Z_0\bigr\}W_1,
  \noenum
  \nexteq{2.5ex}
  W_3=\exp\bigl\{\frac{3}{4}Z_2
                 -\frac{8}{9}Z_1+\frac{17}{36}Z_0\bigr\}W_2,
  \enum
}
其中
\equation{
  Z_i=\eps Z(W_i),\qquad i=0,1,2.
  \enum
}
便得到时刻 $t+\eps$ 处精确解 $V_{t+\eps}$ 的近似。
这些规则使得
\equation{
  W_3=V_{t+\eps}+\rmO(\eps^4)
  \enum
}
因而 $W_3$ 给出所期望的 $V_{t+\eps}$ 的近似。


\lussubsec D.2 夸克流方程的积分

同样有益的是把流方程改写成抽象形式
\equation{
  \partial_t\chi_t=\Delta(V_t)\chi_t,
  \enum
}
其中夸克场 $\chi_t$ 被视为一个复矢量空间的元素。
设在某流时间 $t$ 处给定了 $V_t$ 与 $\chi_t$，
并假定规范场 $W_0,W_1,W_2$ 如上，
则夸克场
\equation{
  \phi_0=\chi_t,
  \noenum
  \nexteq{2.5ex}
  \phi_1=\phi_0+\frac{1}{4}\Delta_0\phi_0,
  \noenum
  \nexteq{2.5ex}
  \phi_2=\phi_0+\frac{8}{9}\Delta_1\phi_1-\frac{2}{9}\Delta_0\phi_0,
  \noenum
  \nexteq{2.5ex}
  \phi_3=\phi_1+\frac{3}{4}\Delta_2\phi_2,
  \enum
}
可以如下计算：
\equation{
  \Delta_i=\eps\Delta(W_i),\qquad i=0,1,2.
  \enum
}
不难证明
\equation{
  \phi_3=\chi_{t+\eps}+\rmO(\eps^4),
  \enum
}
于是 $\phi_3$ 对 $\chi_{t+\eps}$ 的逼近精度
与纯规范流积分器所达到的精度相匹配。



\lusapp E. 伴随流方程 (7.10) 的积分

伴随流方程的数值积分之所以复杂，
是因为规范场流方程的反向积分是指数不稳定的。
本附录先把问题重新表述，
然后讨论其解法。
为简单起见，略去式 (7.10) 中的指标 $k$。


\lussubsec E.1 伴随龙格--库塔积分

对给定步长 $\eps>0$ 及 $t=n\eps$，$n=0,1,2,\ldots$，
令
\equation{
  V_t^{\eps}(x,\mu)\quad\hbox{and}\quad
  \chi_t^{\eps}(x)
  \enum
}
为由附录 D 定义的龙格--库塔积分器生成的规范场与夸克场。
存在另一列夸克场
\equation{
  \xi_s^{\eps}(x),\quad s=t,t-\eps,t-2\eps,\ldots,0,
  \enum
}
其初值为
\equation{
  \xi_t^{\eps}(x)=\eta(x),
  \enum
}
且恒等式
\equation{
  (\xi_s^{\eps},\chi_s^{\eps})=(\eta,\chi_t^{\eps})
  \enum
}
对所有 $s$ 精确成立。
这一性质保证 $\xi_s^{\eps}$ 是伴随流方程的一个近似解，
精度达 $\eps^3$ 阶，
并具有所需的初值。

从 $\chi_s^{\eps}$ 推进到 $\chi_{s+\eps}^{\eps}$
的龙格--库塔步骤相当于把算符
\equation{
  R_s^{\eps}=1+\frac{1}{4}\Delta_0+\frac{3}{4}\Delta_2
  \bigl\{1-\frac{2}{9}\Delta_0+\frac{8}{9}\Delta_1(1+\frac{1}{4}\Delta_0)
  \bigr\}
  \enum
}
作用于 $\chi_s^{\eps}$，
其中 $\Delta_i$ 由式 (D.7) 给出，
而其中出现的场 $W_i$ 由取 $W_0=V_s^{\eps}$ 的式 (D.2) 给出。
由式 (E.4) 即可推知
\equation{
  \xi_s^{\eps}(x)=({R_s^{\eps}}^{\dagger}\xi_{s+\eps}^{\eps})(x)
  \enum
  \nexteq{2.5ex}
  {R_s^{\eps}}^{\dagger}=1+\frac{1}{4}\Delta_0+
  \bigl\{1-\frac{2}{9}\Delta_0+(1+\frac{1}{4}\Delta_0)\frac{8}{9}\Delta_1
  \bigr\}\frac{3}{4}\Delta_2,
  \enum
}
即 $\xi_s^{\eps}$ 通过计算
\equation{
  \lambda_3=\xi_{s+\eps}^{\eps},
  \noenum
  \nexteq{2.5ex}
  \lambda_2=\frac{3}{4}\Delta_2\lambda_3,
  \noenum
  \nexteq{2.5ex}
  \lambda_1=\lambda_3+\frac{8}{9}\Delta_1\lambda_2,
  \noenum
  \nexteq{2.5ex}
  \lambda_0=\lambda_1+\lambda_2+
  \frac{1}{4}\Delta_0(\lambda_1-\frac{8}{9}\lambda_2),
  \enum
}
并由 $\xi_s^{\eps}=\lambda_0$ 得到。


\lussubsec E.2 稳定性问题

伴随递推 (E.8) 是一个光滑化操作，因而在数值上是稳定的。
然而为了能在时刻 $s+\eps$ 应用该递推，
必须已知流时间 $s$ 处的规范场 $V_s^{\eps}$。
原则上，从 $s=t$ 到 $s=0$ 的整列规范场
可以通过梯度流的反向积分重构，
但这样的策略并不实用，
因为沿此方向流是指数不稳定的。

另一种办法是从基本场 $V_0^{\eps}=U$ 出发
通过流的正向积分计算 $V_s^{\eps}$。
若对每个 $s$ 都重新计算一遍，
整个积分所需的规范场更新步数
将按时间步数 $m=t/\eps$ 的平方增长。
该方法在数值上是安全的，
算得的夸克场实际上以机器精度满足恒等式 (E.4)。
但其计算代价可能大得令人却步。


\lussubsec E.3 改进方案

把流时间序列 $0,\eps,\ldots,t$ 分成 $n_b$ 个大小相近的连续区块，
可以使伴随流方程的积分更有效率。
做法是：先用正向积分从时刻 0 算出最后一个区块起点时刻 $r$ 处的场
$V_r^{\eps}$ 并将其保存于计算机内存；
随后从时刻 $t$ 到 $r$ 积分伴随流方程，
其间应用更新步骤 (E.8) 并使用上述安全的规范场重构方法，
只是现在重构从存储的场而非时刻 0 的场开始。

下一步，计算倒数第二个区块中第一个时刻 $r$ 处的规范场并存入内存，
然后伴随流方程的积分照前述方式继续推进到时刻 $r$。
显然各区块可以如此依次处理，
直至积分到达时刻 $0$。

当 $m=n_b^2$ 时该方案的规范场更新总步数最少。
虽然 $m$ 很少恰为整数的平方，
但总可以取 $n_b=\lceil m^{1/2}\rceil$
并把 $m$ 分成大小为 $m_b=\lfloor m/n_b\rfloor$ 与 $m_b+1$ 的区块。
此时计算量近似按 $m^{3/2}$ 标度，
而且在很小的 $m$ 值处就已远小于上一小节安全方案所需的计算量。


\lussubsec E.4 分级方案

进一步加速的办法是把每个区块再分成 $n_b$ 个更小的区块，
后者再分成 $n_b$ 个更小的区块，依此类推，
并在每个区块层级保存一个规范场。
给定 $m$ 与已保存规范场的数目 $n_s$，
接近最优的区块数选取为
\equation{
  n_b=\lceil m^{1/(n_s+1)}\rceil.
  \enum
}
例如对 $m\leq10^3$、$n_s=4$，
需要执行的规范场更新总数不超过约 $5m$。
此外，当多个夸克场同时演化时，
中间规范场只需计算一次。

\begin{thebibliography}{99}



% Wilson flow & emergence of the topological sectors

\bibitem{WilsonFlow}
M. L\"uscher,
{\it Properties and uses of the Wilson flow in lattice QCD},
JHEP 1008 (2010) 071


% Renormalization of the gradient flow

\bibitem{RenFlow}
M. L\"uscher, P. Weisz,
{\it Perturbative analysis of the gradient flow in non-Abelian gauge
theories},
JHEP 1102 (2011) 051


% Step scaling

\bibitem{StepScaling}
M. L\"uscher, P. Weisz, U. Wolff,
{\it A numerical method to compute the running coupling in asymptotically 
free theories},
Nucl. Phys. B359 (1991) 221

\bibitem{NogradiEtAl}
Z. Fodor, K. Holland, J. Kuti, D. Nogradi, C. H. Wong,
{\it The Yang-Mills gradient flow in finite volume},
JHEP 1211 (2012) 007; {\it The gradient flow running coupling scheme},
PoS (Lattice 2012) 050

\bibitem{FritzschRamos}
P. Fritzsch, A. Ramos,
{\it The gradient flow coupling in the Schr\"odinger functional},
arXiv:1301.4388


% Quark source smearing

\bibitem{Guesken}
S. G\"usken,
{\it A study of smearing techniques for hadron correlation functions},
Nucl. Phys. B (Proc. Suppl.) 17 (1990) 361

\bibitem{AlexandrouEtAl}
C. Alexandrou, F. Jegerlehner, S. G\"usken, K. Schilling, R. Sommer,
{\it B meson properties from lattice QCD},
Phys. Lett. B256 (1991) 60


% Renormalization of the Langevin equation

\bibitem{Zinn}
J. Zinn--Justin,
{\it Renormalization and stochastic quantization},
Nucl. Phys. B275 [FS17] (1986) 135

\bibitem{ZinnZwanziger}
J. Zinn--Justin, D. Zwanziger, 
{\it Ward identities for the stochastic quantization of gauge fields},
Nucl. Phys. B295 [FS21] (1988) 297


% Renormalization of quantum field theories on a half space

\bibitem{Symanzik}
K. Symanzik,
{\it Schr\"odinger representation and Casimir effect in
renormalizable quantum field theory},
Nucl. Phys. B190 [FS3] (1981) 1


% Wilson formulation of lattice QCD

\bibitem{Wilson}
K. G. Wilson, {\it Confinement of quarks}, Phys. Rev. D10 (1974) 2445


% Exact chiral symmetry on the lattice

% Ginsparg-Wilson relation

\bibitem{GinspargWilson}
P. H. Ginsparg, K. G. Wilson,
{\it A remnant of chiral symmetry on the lattice},
Phys. Rev. D25 (1982) 2649


% Domain wall fermions

\bibitem{Kaplan}
D. B. Kaplan,
{\it A method for simulating chiral fermions on the lattice},
Phys. Lett. B288 (1992) 342;
{\it Chiral fermions on the lattice},
Nucl. Phys. B (Proc. Suppl.) 30 (1993) 597

\bibitem{Shamir}
Y. Shamir,
{\it Chiral fermions from lattice boundaries},
Nucl. Phys. B406 (1993) 90

\bibitem{FurmanShamir}
V. Furman, Y. Shamir,
{\it Axial symmetries in lattice QCD with Kaplan fermions},
Nucl. Phys. B439 (1995) 54


% Classically perfect Dirac operator

\bibitem{Hasenfratz}
P. Hasenfratz, {\it Prospects for perfect actions},
Nucl. Phys. B (Proc. Suppl.) 63 (1998) 53;
{\it Lattice QCD without tuning, mixing and current renormalization},
Nucl. Phys. B525 (1998) 401

\bibitem{HLN}
P. Hasenfratz, V. Laliena, F. Niedermayer,
{\it The index theorem in QCD with a finite cutoff},
Phys. Lett. B427 (1998) 125


% Neuberger-Dirac operator

\bibitem{NeubergerDirac}
H. Neuberger,
{\it Exactly massless quarks on the lattice},
Phys. Lett. B417 (1998) 141;
{\it More about exactly massless quarks on the lattice},
{\it ibid.}\/ B427 (1998) 353


% Lattice chiral symmetry

\bibitem{ExactChSy}
M. L\"uscher,
{\it Exact chiral symmetry on the lattice and the Ginsparg--Wilson relation},
Phys. Lett. B428 (1998) 342

\bibitem{BoulderReview}
F. Niedermayer,
{\it Exact chiral symmetry, topological charge and related topics},
Nucl. Phys. B (Proc. Suppl.) 73 (1999) 105


% SU(2) chiral perturbation theory

\bibitem{ChPT}
J. Gasser, H. Leutwyler,
{\it Chiral perturbation theory to one loop},
Ann. Phys. 158 (1984) 142


% Global existence of the smoothing flow

\bibitem{ArnoldI}
V. I. Arnold,
{\it Ordinary Differential Equations}, 3rd ed. (Springer-Verlag, Berlin, 2008)


% Stout smearing

\bibitem{Stout}
C. Morningstar, M. Peardon, 
{\it Analytic smearing of SU(3) link variables in lattice QCD},
Phys. Rev. D69 (2004) 054501


% Symanzik effective action

\bibitem{SymanzikSeffI}
K. Symanzik, 
{\it Some topics in quantum field theory},
in: Mathematical problems in theoretical physics, eds. R. Schrader
et al., Lecture Notes in Physics, Vol. 153 (Springer, New York, 1982)

\bibitem{SymanzikSeffII}
K. Symanzik,
{\it Continuum limit and improved action in lattice theories:
(I). Principles and $\varphi^4$ theory},
Nucl. Phys. B226 (1983) 187


% On-shell improvement

\bibitem{OnShell}
M. L\"uscher, P. Weisz,
{\it On-shell improved lattice gauge theories},
Commun. Math. Phys. 97 (1985) 59 [E: {\it ibid.} 98 (1985) 433]


% O(a) improvement

\bibitem{SW}
B. Sheikholeslami, R. Wohlert, 
{\it Improved continuum limit lattice action for QCD with Wilson fermions},
Nucl. Phys. B259 (1985) 572

\bibitem{SFimp}
M. L\"uscher, S. Sint, R. Sommer, P. Weisz,
{\it Chiral symmetry and O(a) improvement in lattice QCD},
Nucl. Phys. B478 (1996) 365


% O(a) improvement with non-degenerate quarks

\bibitem{ImpNonD}
T. Bhattacharya, R. Gupta, W. Lee, S. R. Sharpe, J. M. S. Wu,
{\it Improved bilinears in lattice QCD with nondegenerate quarks},
Phys. Rev. D73 (2006) 034504


% Random sources

\bibitem{MichaelPeisa}
C. Michael, J. Peisa,
{\it Maximal variance reduction for stochastic propagators with 
applications to the static quark spectrum},
Phys. Rev. D58 (1998) 034506


% QCD with open boundary conditions and tm reweighting

\bibitem{TMopenQCD}
M. L\"uscher, S. Schaefer,
{\it  	
Lattice QCD with open boundary conditions and twisted-mass reweighting},
Comput. Phys. Commun. 184 (2013) 519


% Open boundary conditions

\bibitem{openQCD}
M. L\"uscher, S. Schaefer,
{\it Lattice QCD without topology barriers},
JHEP 1107 (2011) 036


% PACS-CS physical point simulations

\bibitem{AokiEtAlI}
S. Aoki et al.~(PACS-CS collab.), 
{\it 2+1 flavor lattice QCD toward the physical point},
Phys. Rev. D79 (2009) 034503

\bibitem{AokiEtAlII}
S. Aoki et al.~(PACS-CS collab.),
{\it Physical point simulation in 2+1 flavor lattice QCD},
Phys. Rev. D81 (2010) 074503


% Iwasaki action

\bibitem{Iwasaki}
Y. Iwasaki,
{\it Renormalization group analysis of lattice theories and
improved lattice action. II -- Four-dimensional non-Abelian SU(N) gauge model},
preprint UTHEP-118 (1983) and arXiv:1111.7054


% Non-perturbative O(a) improvement

\bibitem{AokiEtAlIII}
S. Aoki et al. (CP-PACS collab.),
{\it Nonperturbative O(a) improvement of the Wilson quark action with the 
RG-improved gauge action using the Schr\"odinger functional method},
Phys. Rev. D73 (2006) 034501

\bibitem{KanekoEtAl}
T. Kaneko et al. (CP-PACS/JLQCD and ALPHA collab.),
{\it Non-perturbative improvement of the axial current with 
three dynamical flavors and the Iwasaki gauge action},
JHEP 0704 (2007) 092


% Computation of ZA and ZP

\bibitem{AokiEtAlIV}
S. Aoki et al. (PACS-CS collab.),
{\it Non-perturbative renormalization of quark mass 
in $N_f$=2+1 QCD with the Schr\"odinger functional scheme},
JHEP 1008 (2010) 101


% SF method for ZA

\bibitem{ZASFI}
M. L\"uscher, S. Sint, R. Sommer, H. Wittig,
{\it Non-perturbative determination of the axial current 
normalization constant in O(a) improved lattice QCD},
Nucl. Phys. B491 (1997) 344

\bibitem{ZASFII}
M. Della Morte, R. Hoffmann, F. Knechtli, R. Sommer, U. Wolff,
{\it Non-pertur\-ba\-tive renormalization of the axial current 
with dynamical Wilson fermions},
JHEP 0507 (2005) 007





\end{thebibliography}
